\chapter{Апробация}

\section{Условия апробации}

Апробация проходила в два этапа: внутреннее тестирование на инфраструктуре кафедры и пилотное внедрение в тестовую среду компании-партнёра. Характеристики этапов приведены в таблице~\ref{tab:aprobation_stages}.

\begin{table}[H]
\centering
\caption{Этапы апробации}
\label{tab:aprobation_stages}
\begin{tabular}{|p{4cm}|c|c|p{5.5cm}|c|}
\hline
\textbf{Этап} & \textbf{Период} & \textbf{Длительность} & \textbf{Инфраструктура} & \textbf{Инцидентов} \\
\hline
Лабораторная апробация & сент.--нояб. 2025 & 8 недель & Kubernetes (3 узла), 10 сервисов, Chaos Mesh & 312 \\
\hline
Пилотное внедрение (shadow) & дек. 2025 -- февр. 2026 & 12 недель & Распределённая платформа, 47 сервисов & --- \\
\hline
\end{tabular}
\end{table}

\section{Интеграция с инструментами наблюдаемости}

Апробация подтвердила совместимость системы со стеком наблюдаемости (таблица~\ref{tab:integration_results}).

\begin{table}[H]
\centering
\caption{Результаты интеграционного тестирования}
\label{tab:integration_results}
\begin{tabular}{|p{3.5cm}|p{4cm}|p{5cm}|}
\hline
\textbf{Инструмент} & \textbf{Версия} & \textbf{Результат интеграции} \\
\hline
Prometheus & 2.47 & Сбор 1200 метрик/с без деградации \\
\hline
Elasticsearch & 8.10 & Индексация логов с задержкой $< 2$ с \\
\hline
Jaeger & 1.50 & Анализ трейсов для 47 сервисов \\
\hline
Grafana & 10.2 & Дашборды состояния агентов \\
\hline
OpenTelemetry & 1.28 & Унифицированный сбор телеметрии \\
\hline
Kafka & 3.6 & Шина событий, throughput 5000 msg/s \\
\hline
\end{tabular}
\end{table}

Накладные расходы системы агентов на инфраструктуру: дополнительное потребление CPU --- 8\% от ёмкости кластера, RAM --- 4.2 ГБ на 3 агента, сетевой трафик --- менее 1\% от общего.

\section{Результаты лабораторной апробации}

За 8 недель лабораторного тестирования система продемонстрировала результаты, представленные в таблице~\ref{tab:lab_aprobation_results}.

\begin{table}[H]
\centering
\caption{Результаты лабораторной апробации (312 инцидентов, 8 недель)}
\label{tab:lab_aprobation_results}
\begin{tabular}{|p{5cm}|c|c|}
\hline
\textbf{Показатель} & \textbf{Значение} & \textbf{Целевое} \\
\hline
Обнаружено инцидентов & 298 из 312 (95.5\%) & $> 90\%$ \\
\hline
Корректно диагностирована причина & 271 из 298 (90.9\%) & $> 85\%$ \\
\hline
Успешное автовосстановление & 243 из 271 (89.7\%) & $> 80\%$ \\
\hline
Среднее время обнаружения & 2.3 мин & $< 5$ мин \\
\hline
Среднее время диагностики & 4.7 мин & $< 10$ мин \\
\hline
Ложные срабатывания & 24 (7.7\%) & $< 15\%$ \\
\hline
Некорректные действия Safe Executor & 0 & 0 \\
\hline
\end{tabular}
\end{table}

Классификация ошибок системы представлена в таблице~\ref{tab:error_breakdown}.

\begin{table}[H]
\centering
\caption{Классификация ошибок системы (лабораторная апробация)}
\label{tab:error_breakdown}
\begin{tabular}{|p{3.5cm}|c|p{6cm}|}
\hline
\textbf{Тип ошибки} & \textbf{Кол-во} & \textbf{Причина} \\
\hline
\multicolumn{3}{|l|}{\textit{Пропущенные инциденты (14)}} \\
\hline
Медленные утечки ресурсов & 9 & Период развития $> 4$ часов \\
\hline
Нетипичные метрики & 5 & Отсутствовали в обучающей выборке \\
\hline
\multicolumn{3}{|l|}{\textit{Некорректная диагностика причины (28)}} \\
\hline
Каскадные отказы & 28 & Затронуто $> 5$ сервисов одновременно \\
\hline
\end{tabular}
\end{table}

\section{Результаты пилотного внедрения}

В режиме shadow-мониторинга система анализировала инциденты параллельно с дежурными инженерами. Сравнение проводилось по двум критериям: скорость определения корневой причины и корректность диагноза. Результаты представлены в таблице~\ref{tab:shadow_comparison}.

\begin{table}[H]
\centering
\caption{Сравнение: система агентов и дежурные инженеры (shadow-режим, 12 недель)}
\label{tab:shadow_comparison}
\begin{tabular}{|p{4.5cm}|c|c|c|}
\hline
\textbf{Метрика} & \textbf{Инженеры} & \textbf{Агенты} & \textbf{$\Delta$} \\
\hline
Среднее время диагностики, мин & $38.4 \pm 15.2$ & $6.1 \pm 2.8$ & $-84\%$ \\
\hline
Корректность диагноза, \% & 87.3 & 89.1 & $+1.8$ п.п. \\
\hline
Пропущенные инциденты, \% & 4.2 & 5.8 & $+1.6$ п.п. \\
\hline
Ложные тревоги в сутки & $1.2$ & $2.8$ & $+133\%$ \\
\hline
\end{tabular}
\end{table}

Система агентов значительно превосходит инженеров по скорости диагностики ($-84\%$) при сопоставимой корректности. Инженеры реже пропускают инциденты (4.2\% против 5.8\%), но это преимущество нивелируется при комбинированном использовании: совместная работа агентов и инженеров снизила долю пропущенных инцидентов до 1.4\%.

Количество ложных тревог у агентов выше (2.8 против 1.2 в сутки), что требует дальнейшей оптимизации. Анализ показал, что 60\% ложных тревог связаны с плановыми деплоями, которые система некорректно классифицировала как аномалии.

\section{Оценка безопасности}

За весь период апробации (20 недель) система Safe Executor:

\begin{itemize}
    \item Заблокировала 18 потенциально опасных действий, предложенных модулем восстановления (перезапуск критических сервисов при высокой нагрузке, масштабирование сверх лимитов ресурсов);
    \item Не допустила ни одного некорректного действия восстановления, прошедшего валидацию;
    \item Средняя задержка валидации: 1.2 с (при пороге $r_{\max} = 0.3$).
\end{itemize}
